% Source: https://artofproblemsolving.com/community/c3h1993958_2001_national_countdown_problem
% Stable identity: aops-topic-1993958-q01
A rectangle with perimeter $144$ has length and width $\ell$ and $72-\ell$, where $1\le\ell\le71$ is an integer. Its area is
\[
\ell(72-\ell)=1296-(36-\ell)^2 .
\]
The area is largest, $1296$, when $\ell=36$ (the $36\times36$ square), and smallest when $|36-\ell|$ is as large as possible, $\ell=1$ or $71$, giving $1\cdot71=71$.

\begin{asy}
usepackage("amsmath");
size(7cm);
defaultpen(fontsize(10pt));
draw((0,0)--(1.8,0)--(1.8,1.8)--(0,1.8)--cycle,linewidth(0.9));
label("$36$",(0.9,0),S);label("$36$",(1.8,0.9),E);
label("area $1296$",(0.9,0.9));
draw((2.6,0.84)--(6.15,0.84)--(6.15,0.96)--(2.6,0.96)--cycle,linewidth(0.9));
label("$71$",(4.375,0.84),S);label("$1$",(6.15,0.9),E);
label("area $71$",(4.375,0.96),N);
label("\mbox{}",(min(currentpicture,user=true)+max(currentpicture,user=true))/2,black); // keeps the text after the figure black when LaTeX typesets the labels inline
\end{asy}

The greatest possible difference is $1296-71=\boxed{1225}\text{ cm}^2$.
